\documentclass[10pt]{article}
\usepackage[a4paper,margin=0.6in]{geometry}
\usepackage{amsmath,amssymb,amsfonts,mathtools}
\usepackage[dvipsnames]{xcolor}
\usepackage{array,booktabs,float,longtable}
\usepackage[colorlinks=true,linkcolor=blue,citecolor=blue]{hyperref}
\allowdisplaybreaks
\numberwithin{equation}{section}
\setlength{\jot}{4pt}
\setlength{\emergencystretch}{3em}
\newcommand{\new}[1]{{\color{blue}#1}}
\newcommand{\old}[1]{{\color{red}#1}}
\newcommand{\Lv}{L_{\vee}}
\newcommand{\Ls}{L_{*}}
\newcommand{\OO}{\mathcal O}
\newcommand{\KK}{\mathcal K}
\newcommand{\JJ}{\mathcal J}
\newcommand{\BB}{\mathcal B}
\newcommand{\NN}{\mathcal N}
\newcommand{\PP}{\mathcal P}
\newcommand{\LL}{\mathcal L}
\newcommand{\Eh}{\hat E}
\newcommand{\Dh}{\hat\Delta}
\newcommand{\pre}{\frac{1}{(2\pi)^3}}
\newcommand{\phw}{e^{-ik\cdot(w'-w)}}
\newcommand{\Tr}{{\rm Tr}}
\newcommand{\FS}[1]{(FS~#1)}
\newcommand{\EC}[1]{(EC~#1)}

\title{The $\log^2$ term, and the final finite result in colour-singlet form}
\date{}

\begin{document}
\maketitle
\tableofcontents

%==========================================================================================
\section{Summary}\label{sec:summary}
References: ``FS~$n$'' is Eq.~($n$) of \texttt{Finite\_terms\_simplified.pdf} (updated with this note); ``EC~$n$'' is Eq.~($n$) of \texttt{Evolution\_terms\_corrected.pdf}. Notation as in FS: $a=\frac{k^+}{\vee}$, $\Lv=\log\frac{\vee}{k^+}$, $\Ls=\log\frac{\sqrt{\vee\Lambda}}{k^+}$, $\NN=\pre\frac{g^4}{16\pi^5k^+}$.

\paragraph{1. The $\log^2$ term (Sec.~\ref{sec:log2}).} It is in the final result, and that is correct for your file as it stands.
\begin{itemize}
\item It appears once in \texttt{Finite\_terms\_1.tex}: in the bracket of Group I Row II part 1, $-\frac{67}{9}+\frac{\pi^2}{3}+\log^2(1-a)+\log(1-a)\Lv$. No other row has a $\log^2$, so nothing cancels it.
\item It is the finite part of the dressed one-$\rho$ state $|\Psi^A_{g\rho}\rangle+|\Psi^B_{g\rho}\rangle$. I re-derived it analytically (Sec.~\ref{sec:origin}): it comes from the lower limit $\xi\ge\frac{k^+}{\vee}$ of the $\xi$ integral of $|\Psi^B_{g\rho}\rangle$, i.e.\ from requiring that the parent gluon has momentum $\frac{k^+}{\xi}<\vee$.
\item The same limit gives the $\log(1-a)$ in the pole bracket of Row II part 1, i.e.\ the uncancelled pole of Q1. The $\log^2$ and the Q1 pole come together.
\begin{itemize}
\item If the limit is kept (as in your file), both stay.
\item If it is dropped where it is not needed, both disappear and Row II part 1 leaves only $\frac{67}{9}-\frac{\pi^2}{3}$. ``Not needed'' means the parts of the integrand that are regular at $\xi\to0$; this is AABKL's $\int_0^1d\xi$.
\item One cannot drop the $\log^2$ and keep the rest.
\end{itemize}
\item There is no separate $\log^2$ in the result: together with its partner,
\begin{equation*}
-\log^2(1-a)-\Lv\log(1-a)=-\log\Big(1-\frac{k^+}{\vee}\Big)\log\Big(\frac{\vee}{k^+}-1\Big),
\end{equation*}
which is $\approx a\log\frac1a\to0$ for $k^+\ll\vee$.
\item It is not a running-coupling term. In the final result it is now written separately as $V_\vee$, next to the Q1 pole, and the running-coupling remainder is the CMW constant $\frac{67}{9}-\frac{\pi^2}{3}$ alone. FS~(4.3) and FS~(11.1) are updated accordingly.
\end{itemize}

\paragraph{2. Final result in colour-singlet form (Secs.~\ref{sec:rules}--\ref{sec:final}).} The projectile is in a colour singlet,
\begin{equation*}
\langle\rho^a(u)\rho^b(v)\rangle=\delta^{ab}\,\frac{\langle\rho^f(u)\rho^f(v)\rangle}{N_c^2-1}\equiv\delta^{ab}\,\PP(u,v).
\end{equation*}
\begin{itemize}
\item Every colour factor of the final result FS~(11.1) is projected and written with traces of adjoint $U$, $U^\dagger$, $T^e$.
\item This covers LO, every LO-like term, the four two-$\rho$ rows, and the two-$\rho$ parts of all three-$\rho$ rows (merge by merge).
\item The index placements are those valid for a general adjoint $U$ (Sec.~\ref{sec:placement}). For $U=U^T$ (your convention $U^{bc}=U^{cb}$) they reduce to your colour factors.
\item The three-$\rho$ symmetric parts need $\langle\rho\rho\rho\rangle$ and vanish for a Gaussian projectile. The four-$\rho$ term (Q2) needs $\langle\rho\rho\rho\rho\rangle$.
\item The result is Eq.~\eqref{eq:final}.
\end{itemize}
\paragraph{\new{3. Correction: $H_4$ (Eq.~\eqref{eq:sH4}).}} \new{The DGLAP integral, Eq.~\eqref{eq:sDG}, has the lower limit $\xi=\frac{k^+}{\vee}$. Its small-$\xi$ end (parent gluon $\frac{k^+}{\xi}\gg k^+$) contains $\log\frac{\vee}{k^+}H_4$, a BFKL term that is already in $\BB_s$. The earlier version of Eq.~\eqref{eq:final} kept it in the DGLAP term as well, i.e.\ counted $\frac12\log\frac\vee\Lambda H_4$ twice. Eq.~\eqref{eq:final} now has $-\Lv H_4$ in the DGLAP line (details: \texttt{BFKL\_terms\_check.pdf}, Sec.~8).}

\textbf{Checks} (Sec.~\ref{sec:checks}): every trace form equals the direct contraction of its colour network with random general SU(3) Wilson lines, to @@TRMAX@@. The general-$U$ colour operators equal the first-principles colour functions, to @@FPMAX@@.

%==========================================================================================
\section{The \texorpdfstring{$\log^2$}{log2} term}\label{sec:log2}
\subsection{Where it is}
Group I Row II part 1 (the first equation of Group I Row II in \texttt{Finite\_terms\_1.tex}), bracket
\begin{equation}
\Big[-\frac2\epsilon-2\gamma_E-\log\frac{(x-w)^2\mu^2_{\overline{MS}}}{4}\Big]\Big[\frac{11}{3}+\log\frac{\vee}{k^+}+\log\Big(\frac{\vee}{k^+}-1\Big)\Big]-\frac{67}{9}+\frac{\pi^2}{3}+\log^2\Big(1-\frac{k^+}{\vee}\Big)+\log\Big(1-\frac{k^+}{\vee}\Big)\log\frac{\vee}{k^+}.
\label{eq:bracket}
\end{equation}
With its c.c.\ it gives FS~(4.1). The $\beta_0$ term cancels against Row V 2b UV. The rest is
\begin{equation}
\frac{\alpha_s}{2\pi}N_c\Big[\frac{67}{9}-\frac{\pi^2}{3}\Big]{\rm LO}\;\underbrace{-\;\frac{\alpha_s}{2\pi}N_c\Big[\log^2(1-a)+\Lv\log(1-a)\Big]{\rm LO}}_{V_\vee}
\;+\;\text{pole terms}\ \propto\big(2\Lv+\log(1-a)\big).
\end{equation}

\subsection{Where it comes from}\label{sec:origin}
\texttt{Calculation\_main.tex}, Eqs.~for $|\Psi^A_{g\rho}\rangle$ and $|\Psi^B_{g\rho}\rangle$ (one $\rho$, common factor $\frac{g^3N_c\rho^a(-k)k^i}{32\pi^{7/2}k^2}$, $P_\epsilon\equiv-\frac2\epsilon+\log\frac{k^2}{\mu^2_{\overline{MS}}}$):
\begin{align*}
|\Psi^A_{g\rho}\rangle&:\quad P_\epsilon\Big[\frac{11}{3}+\log\frac{\Lambda}{k^+}-2\log\frac{\vee}{\Lambda}\Big]+\frac{\pi^2}{3}-\frac{67}{9}+\log^2\frac{\Lambda}{k^+}-\log^2\frac{\vee}{k^+},\\
|\Psi^B_{g\rho}\rangle&:\quad\int_\Lambda^\vee\frac{dk^+}{\sqrt{k^+}}\int_{\Lambda/k^+}^{1-\Lambda/k^+}\frac{d\xi}{\sqrt\xi(1-\xi)}\Big[(\xi-2)\log\xi+(\xi+1)\log(1-\xi)-(\xi-2)P_\epsilon\Big]\,|g(\xi k^+)\rangle.
\end{align*}
In $|\Psi^B\rangle$ the gluon has momentum $\xi k^+$; $k^+$ is the momentum of its parent. Call the gluon momentum $k^+$ again, i.e.\ substitute $k^+\to k^+/\xi$ (Jacobian included, the measure becomes $\frac{d\xi}{\xi(1-\xi)}$). The limits become:
\begin{itemize}
\item parent below the cutoff, $\frac{k^+}{\xi}<\vee$: \quad $\xi>\frac{k^+}{\vee}=a$;
\item $1-\xi>\frac{\Lambda\xi}{k^+}$: \quad $\xi<1-\delta+O(\delta^2)$, $\delta=\frac\Lambda{k^+}$.
\end{itemize}
So
\begin{equation}
|\Psi^B\rangle\ \to\ \int_a^{1-\delta}\frac{d\xi}{\xi(1-\xi)}\Big[(\xi-2)\log\xi+(\xi+1)\log(1-\xi)-(\xi-2)P_\epsilon\Big].
\label{eq:psiB}
\end{equation}
\textbf{Pole part:} $\frac{2-\xi}{\xi(1-\xi)}=\frac2\xi+\frac1{1-\xi}$, so
\begin{equation*}
\int_a^{1-\delta}\Big(\frac2\xi+\frac{1}{1-\xi}\Big)=2\Lv+\log(1-a)+\log\frac{k^+}{\Lambda}.
\end{equation*}
\textbf{Finite part}, term by term:
\begin{center}\small
\begin{tabular}{@{}lll@{}}\toprule
integrand & $\int_a^{1-\delta}d\xi$ & comes from\\\midrule
$-2\,\frac{\log\xi}{\xi}$ & $\Lv^2$ & soft end $\xi\to0$ (needs the limit $a$)\\
$-\frac{\log\xi}{1-\xi}$ & ${\rm Li}_2(1-a)$ & \\
$\frac{\log(1-\xi)}{\xi}$ & ${\rm Li}_2(a)-\frac{\pi^2}{6}$ & \\
$2\,\frac{\log(1-\xi)}{1-\xi}$ & $\log^2(1-a)-\log^2\frac{\Lambda}{k^+}$ & \textbf{the $\log^2(1-a)$}\\\bottomrule
\end{tabular}
\end{center}
With ${\rm Li}_2(a)+{\rm Li}_2(1-a)=\frac{\pi^2}{6}-\log a\log(1-a)$ the sum is $\Lv^2+\Lv\log(1-a)+\log^2(1-a)-\log^2\frac{\Lambda}{k^+}$. Adding $|\Psi^A\rangle$:
\begin{equation}
P_\epsilon\Big[\frac{11}{3}+2\Lv+\log(1-a)-2\log\frac{\vee}{\Lambda}\Big]+\frac{\pi^2}{3}-\frac{67}{9}+\log^2(1-a)+\Lv\log(1-a).
\end{equation}
This is exactly the bracket of Eq.~\eqref{eq:bracket} (your $\log\frac{\vee}{k^+}+\log(\frac{\vee}{k^+}-1)=2\Lv+\log(1-a)$, and the $-2\log\frac\vee\Lambda$ is the evolution part). Numerically (\texttt{log2\_origin.py}), $\Psi^A+\Psi^B$ minus the bracket is @@L2FIN@@ (finite part) and @@L2POL@@ (pole part), i.e.\ $O(\delta\log\delta)$.

\subsection{What this means}
\begin{enumerate}
\item \textbf{It is correct for your file}, and it is in the final result. Nothing cancels it:
\begin{itemize}
\item it is the only $\log^2$ in \texttt{Finite\_terms\_1.tex};
\item it multiplies the LO structure. The only other terms with that structure are Row V 2b UV (no $\log(1-a)$) and Group II Row IV. Group II Row IV, in the plus distribution restricted to $\xi>a$, has a single $\log(1-a)$, which is absorbed in $D_{h/g}$ (FS~(5.1)).
\end{itemize}
\item \textbf{Compact form.} Exactly,
\begin{equation}
V_\vee\equiv-\log^2(1-a)-\Lv\log(1-a)=-\log(1-a)\log\frac{1-a}{a}=-\log\Big(1-\frac{k^+}{\vee}\Big)\log\Big(\frac{\vee}{k^+}-1\Big).
\label{eq:Vv}
\end{equation}
\begin{itemize}
\item It is small unless $k^+$ is close to $\vee$: $V_\vee/(a\log\frac1a)=$ @@SMALLA@@ for $a=10^{-2},10^{-4},10^{-6}$.
\item It grows only at the end point $k^+\to\vee$, like $-\log^2(1-a)$.
\end{itemize}
\item \textbf{It comes with the Q1 pole.} The $\log(1-a)$ of the pole, the $\log^2(1-a)$ and the $\Lv\log(1-a)$ all come from the limit $\xi>a$ applied to the four terms of Eq.~\eqref{eq:psiB} that are regular at $\xi\to0$: $\frac{1}{1-\xi}$, $\frac{\log\xi}{1-\xi}$, $\frac{\log(1-\xi)}{\xi}$, $\frac{\log(1-\xi)}{1-\xi}$.
\begin{itemize}
\item The terms $\frac2\xi$ and $-\frac{2\log\xi}{\xi}$ do need the limit: they give $2\Lv$ and $\Lv^2$, the rapidity logarithms.
\item If the four regular terms are integrated from $\xi=0$ (as in AABKL, where the virtual term is $\int_0^1d\xi$), the pole bracket is $\frac{11}{3}+2\Lv-2\log\frac\vee\Lambda$ and the finite part is $\frac{\pi^2}{3}-\frac{67}{9}$ only. Numerically: @@L2UNR@@.
\item Then the pole tally of FS Sec.~5 closes ($\frac{11}{6}$, absorbed in $D_{h/g}$), Q1 disappears, and so does $V_\vee$.
\end{itemize}
\item \textbf{So there are two consistent options:}
\begin{itemize}
\item[(i)] keep the limit $\xi\ge\frac{k^+}{\vee}$ everywhere (your file): $V_\vee$ and the Q1 pole are both in the result;
\item[(ii)] apply it only where it is needed (the soft end): neither is.
\end{itemize}
Which one is right is the question Q1. The final result below is written for (i), with $V_\vee$ and the Q1 term marked; for (ii) delete both lines.
\end{enumerate}

%==========================================================================================
\section{Colour-singlet projection: rules}\label{sec:rules}
\subsection{Projection and traces}
\begin{itemize}
\item $\langle\rho^a(u)\rho^b(v)\rangle=\delta^{ab}\PP(u,v)$, $\PP(u,v)\equiv\frac{\langle\rho^f(u)\rho^f(v)\rangle}{N_c^2-1}$. $\PP$ is symmetric (the commutator term is $\propto\langle\rho^c\rangle=0$), so the order of the two charges does not matter.
\item After the projection, all colour indices are contracted among Wilson lines and structure constants. The network is written with traces, as in \texttt{Evolution\_singlet.pdf}:
\begin{itemize}
\item $(T^e)_{bc}=-if^{ebc}$; $U$ is real orthogonal, so $U^T=U^\dagger$;
\item two $f$'s joined by a chain $C$ of Wilson lines: $\sum_dC^{ad}T^d=C^\dagger T^aC$, giving ${\rm Tr}\{T^eAT^eB\}$ or ${\rm Tr}\{AT^e\}{\rm Tr}\{BT^e\}$;
\item $UU^\dagger=1$, ${\rm Tr}\,1=N_c^2-1$, $T^eT^e=N_c$, ${\rm Tr}\{T^eAT^eA^\dagger\}=N_c{\rm Tr}A$.
\end{itemize}
\item After the projection, terms that differ only in the point of one Wilson line are combined into differences, e.g.\ ${\rm Tr}\{(U(x)-U(y))(U^\dagger(x')-U^\dagger(y'))\}$.
\item Every trace below is real (for real $U$; a trace with two $T^e$ is real too). So the c.c.\ of a Group I or Group III term is: complex-conjugate every number (also the $i$'s of the colour factors), then exchange $w\leftrightarrow w'$ (resp.\ $w\leftrightarrow y'$), as in FS Sec.~2.1.
\end{itemize}

\subsection{Index placement for a general \texorpdfstring{$U$}{U}}\label{sec:placement}
Your files use $U^{bc}=U^{cb}$. For the projection the order matters: for a general $U$, a misplaced index gives, e.g., ${\rm Tr}\{U(p)U(q)\}$ instead of ${\rm Tr}\{U^\dagger(p)U(q)\}$. The projections below use:
\begin{itemize}
\item \textbf{Group I} (Rows II, III, IV, V): your colour factors as written. They agree with first principles for a general $U$ (\texttt{Group\_I\_check.pdf}); re-checked here: @@G1MAX@@.
\item \textbf{LO}: $[U(y')-U(x')]^{ce'}\rho^{e'}(x')\,[U(y)-U(x)]^{ce}\rho^e(x)$. Your $U^{ce}(x)$ next to $U^{ec}(y)$ is the same for $U=U^T$.
\item \textbf{Groups II and III}: the placements of correction M4 (\texttt{Groups\_II\_III\_check.pdf}, \texttt{fixed23.py}). Every changed Wilson line is a $U(z)$ or $U(y')$:
\begin{center}\footnotesize
\begin{tabular}{@{}p{0.42\textwidth}ll@{}}\toprule
where & your file & general $U$\\\midrule
$\mathcal R^{ab}(w,z,x)$ (all of Group II Row I, Group III Rows I, II, III(2)) & $f^{acd}U^{bc}(z)U^{de}(x)\rho^e(x)$ & $f^{acd}U^{cb}(z)U^{de}(x)\rho^e(x)$\\
Group II Row I, two-charge bracket & $U^{bc'}(z)U^{ad'}(x')U^{c'e'}(y)$ & $U^{c'b}(z)U^{ad'}(x')U^{c'e'}(y)$\\
Group III Row I & $U^{bd'}(z)U^{d'e'}(y)U^{ac'}(w')$ & $U^{d'b}(z)U^{d'e'}(y)U^{ac'}(w')$\\
Group III Row II & $U^{c'd'}(y)U^{ab'}(x')U^{bc'}(z)$ & $U^{c'd'}(y)U^{ab'}(x')U^{c'b}(z)$\\
Group III Rows III, V, VI, III(2) (conjugate side) & $U^{e'c'}(y')$, $U^{bd'}(z)U^{c'e'}(x')$ & $U^{c'e'}(y')$, $U^{d'b}(z)U^{c'e'}(x')$\\
Group III Rows III(1), VI (amplitude side) & $U^{bc}(z)U^{ad}(y)U^{ce}(x)$ & $U^{cb}(z)U^{ad}(y)U^{ce}(x)$\\
Group II Row I part 2, $\OO_{58}$ & see Eq.~\eqref{eq:O58} & \\\bottomrule
\end{tabular}
\end{center}
\item With these placements every colour operator equals the first-principles colour function of \texttt{fixed23.py} for a general $U$ (@@GSMAX@@). For $U=U^T$ it equals your operator (@@GSSYM@@). As written, your operators differ for a general $U$ (by 14\%--220\%).
\item $\OO_{58}$ with the general placement is the product of the two $\mathcal R$'s:
\begin{align}
\OO_{58}=\mathcal R^{ab}(w',z,x')\,\mathcal R^{ab}(w,z,x)={}&f^{c'bd'}f^{cbd}U^{ac'}(w')U^{ac}(w)\rho^{d'}(x')\rho^d(x)\notag\\
&-f^{c'bd'}f^{acd}U^{ac'}(w')U^{cb}(z)U^{de}(x)\rho^{d'}(x')\rho^e(x)\notag\\
&-f^{ac'd'}f^{cbd}U^{ac}(w)U^{c'b}(z)U^{d'e'}(x')\rho^{e'}(x')\rho^d(x)\notag\\
&+N_cU^{de'}(x')U^{de}(x)\rho^{e'}(x')\rho^e(x).
\label{eq:O58}
\end{align}
\end{itemize}

\subsection{Three and four \texorpdfstring{$\rho$}{rho}'s}
\begin{itemize}
\item A three-$\rho$ colour operator is $\OO=S(\OO)+\OO^{(2)}$ (FS~(9.1)). $\OO^{(2)}$ has two $\rho$'s and is projected like any two-$\rho$ term (Sec.~\ref{sec:three}).
\item $S(\OO)$ needs the three-point function. For a colour singlet, the only invariant tensors are $d^{abc}$ and $f^{abc}$:
\begin{equation}
\langle S\big(\rho^a(u)\rho^b(v)\rho^c(t)\big)\rangle=\frac{N_c\,d^{abc}}{(N_c^2-4)(N_c^2-1)}\langle d^{efg}\rho^e(u)\rho^f(v)\rho^g(t)\rangle+\frac{f^{abc}}{N_c(N_c^2-1)}\langle f^{efg}\,S\big(\rho^e(u)\rho^f(v)\rho^g(t)\big)\rangle.
\end{equation}
Both vanish for a Gaussian (McLerran--Venugopalan-type) projectile.
\item Four $\rho$'s (FS~(10.1), Q2): $\langle\rho\rho\rho\rho\rangle$; for a Gaussian projectile it is the sum of the three pairings, $\PP\PP$. The term is zero if Q2 is answered with the range $(\Lambda,\vee)$.
\end{itemize}

%==========================================================================================
\section{LO and the LO-like terms}\label{sec:LO}
\begin{equation}
\boxed{\LL[f]\equiv\pre\frac{g^2}{2\pi^2k^+}\int_{x,y,x',y'}e^{-ik\cdot(y'-y)}\,f\,\frac{(x-y)\cdot(x'-y')}{(x-y)^2(x'-y')^2}\,\PP(x,x')\,{\rm Tr}\big\{\big(U^\dagger(x)-U^\dagger(y)\big)\big(U(x')-U(y')\big)\big\}}
\label{eq:LO}
\end{equation}
($x$, $x'$: charges; $y$, $y'$: the measured gluon in the amplitude and its conjugate). $\LL_\xi[f]$: the same with $e^{-ik\cdot(y'-y)}\to e^{-ik\cdot(y'-y)/\xi}$. The projected LO colour factor is @@LOTR@@ (@@LOCHK@@).

The LO-like terms of FS~(11.1) in singlet form ($X=x-y$, $X'=x'-y'$, $\Delta=y'-y$):
\begin{align}
\frac{d^3N}{d^3k}\Big|_{\rm RC}&=\frac{\alpha_s}{2\pi}N_c\Big[\frac{67}{9}-\frac{\pi^2}{3}\Big]\LL[1],\label{eq:sRC}\\
\frac{d^3N}{d^3k}\Big|_{V_\vee}&=-\frac{\alpha_s}{2\pi}N_c\log\Big(1-\frac{k^+}{\vee}\Big)\log\Big(\frac{\vee}{k^+}-1\Big)\LL[1]\qquad\text{(with Q1)},\label{eq:sVv}\\
\frac{d^3N}{d^3k}\Big|_{\rm Q1}&=\frac{\alpha_sN_c}{4\pi}\log\Big(1-\frac{k^+}{\vee}\Big)\LL\big[\Eh_X+\Eh_{X'}\big],\label{eq:sQ1}\\
\frac{d^3N}{d^3k}\Big|_{\Ls}&=\frac{\alpha_sN_c}{2\pi}\,\Ls\,\LL\Big[\log\frac{(y-y')^4}{(x-y)^2(x'-y')^2}\Big],\label{eq:sLs}\\
\frac{d^3N}{d^3k}\Big|_{\rm DGLAP}&=-\frac{\alpha_s}{2\pi}\int_{a}^1\frac{d\xi}{\xi^2}P^{(r)}_{gg,+}(\xi)\,\LL_\xi\big[\ell_\Delta(\mu_F)\big],\label{eq:sDG}\\
\new{H_4}&\new{{}=2N_c\,\NN\int_{x,x',z,w,w'}\phw\,\KK(w-z)\cdot\KK(w'-z)\;\KK(x-z)\cdot\KK(x'-z)}\notag\\
&\qquad\new{\times\PP(x,x')\,{\rm Tr}\big\{\big(U^\dagger(z)-U^\dagger(x')\big)\big(U(z)-U(x)\big)\big\},}\label{eq:sH4}\\
\frac{d^3N_{\rm lin}}{d^3k}&=\frac{\alpha_s}{2\pi}\int_{a}^1\frac{d\xi}{\xi^2}2N_c\xi(1-\xi)\pre\frac{g^2}{2\pi^2k^+}\int_{x,y,x',y'}e^{-ik\cdot\Delta/\xi}\Big[\KK(X)\cdot\KK(X')-2\big(\KK(X)\cdot\Dh\big)\big(\KK(X')\cdot\Dh\big)\Big]\notag\\
&\qquad\times\PP(x,x')\,{\rm Tr}\big\{\big(U^\dagger(x)-U^\dagger(y)\big)\big(U(x')-U(y')\big)\big\},\label{eq:slin}
\end{align}
with $\Eh_r=\frac2\epsilon+\log\frac{r^2\mu^2_{\overline{MS}}}{c_0^2}$, $\ell_\Delta=\log\frac{\Delta^2\mu_F^2}{c_0^2}$, $P^{(r)}_{gg,+}=2N_c[\frac{\xi}{(1-\xi)_+}+\frac{1-\xi}{\xi}+\xi(1-\xi)]$, $\Dh=\Delta/|\Delta|$. Eq.~\eqref{eq:sLs} is one of the $\Ls$ terms that add up to the remainder of Eq.~\eqref{eq:final}. \new{So is $\Ls H_4$: Eq.~\eqref{eq:sDG} $=\Lv H_4+(\text{finite for }\frac{k^+}{\vee}\to0)$, with $\Lv H_4=\Ls H_4+\frac12\log\frac\vee\Lambda H_4$, and both parts are already in the evolution line of Eq.~\eqref{eq:final}.}

%==========================================================================================
\section{The two-\texorpdfstring{$\rho$}{rho} rows}\label{sec:two}
The kernels are those of FS Sec.~7, unchanged; only the colour operator is replaced by $\PP(u,v)\,C$, with $u,v$ the two charges.

\paragraph{Group I Row V, 2a} (FS~(7.6), $+$c.c.):
\begin{align}
\frac{d^3N_{2a}}{d^3k}&=\pre\frac{g^4}{16\pi^5k^+}\int_{x,y,z,x',w'}\{\dots\}_{\rm V2}\,\PP(x',y)\,C_{2a}+{\rm c.c.},\qquad\{\dots\}_{\rm V2}\ \text{of FS~(7.5)},\label{eq:s2a}\\
C_{2a}&=@@V2a@@.\notag
\end{align}
\paragraph{Group I Row V, 2b finite} (FS~(7.7), $+$c.c.):
\begin{align}
\frac{d^3N_{2b,\rm fin}}{d^3k}&=\pre\frac{g^4}{16\pi^5k^+}\int_{x,y,z,x',w'}\Big\{\{\dots\}_{\rm V2}\,\PP(x',y)\,C_{2b}\notag\\
&\qquad+e^{-ik\cdot(w'-z)}\Big(\frac{11}{6}+2\Ls\Big)\frac{\KK(x'-w')\cdot\KK(y-z)\;(y-x)\cdot(y-z)}{(y-x)^2(x-z)^2}\,\PP(x',y)\,C_{2b}'\Big\}+{\rm c.c.},\label{eq:s2b}\\
C_{2b}&=@@V2b1@@,\notag\\
C_{2b}'&=@@V2b2@@.\notag
\end{align}
\paragraph{Group II Row I part 2} (FS~(7.9), no c.c.; $c_\lambda$, $c_{\lambda'}$, $c_L$, $c_\ell$ of FS~(7.9)):
\begin{align}
\frac{d^3N_{\rm I(2)}}{d^3k}&=\pre\frac{g^4}{8\pi^5k^+}\int_{z,w,w',x,x'}\phw\Big[c_\lambda\lambda+c_{\lambda'}\lambda'+c_L\Ls+c_\ell\log(1-a)\notag\\
&\qquad-\frac{d_\perp}{4}\frac{1-a}{(z-w)^2(z-w')^2}\Big]\PP(x',x)\,C_{58},\label{eq:s58}\\
C_{58}&=@@G2I2@@.\notag
\end{align}
\paragraph{Group III Row III part 2} (FS~(7.12), $+$c.c.; $c_{1..5}$, $\mathbb I_{1..5}$, $s$, $\mathcal G$, $\mathcal J(b)$ of FS Sec.~7.3):
\begin{align}
\frac{d^3N_{\rm III(2)}}{d^3k}&=\pre\frac{g^4}{8\pi^5}\int_{w,x',y',z,x}e^{-ik\cdot(y'-w)}\Big\{c_1\mathbb I_1+\frac{c_2}{k^+}\mathbb I_2+c_3\mathbb I_3+\frac{c_4}{k^{+2}}\mathbb I_4+\frac{c_5}{k^+}\mathbb I_5\notag\\
&\qquad-\frac{s}{k^+B}\big[\mathcal G+\Ls+\mathcal J(b)\big]\Big\}\PP(x',x)\,C_{\rm III(2)}+{\rm c.c.},\label{eq:sIII2}\\
C_{\rm III(2)}&=@@III2@@.\notag
\end{align}

%==========================================================================================
\section{The three-\texorpdfstring{$\rho$}{rho} rows: two-\texorpdfstring{$\rho$}{rho} parts}\label{sec:three}
For a three-$\rho$ row with charges $(u_1,u_2,u_3)$ and kernel $\KK_{\rm row}$, the two-$\rho$ part is
\begin{equation}
\frac{d^3N^{(2)}_{\rm row}}{d^3k}=\pre\int(\text{phase})\sum_{\rm merges}\KK_{\rm row}\big|_{u_j\to u_i}\;\PP(u_i,u_k)\;C^{\rm row}_{u_j\to u_i}\;(+{\rm c.c.}),
\label{eq:three}
\end{equation}
where $u_j\to u_i$ is the merged pair (the point $u_j$ is set equal to $u_i$; the integral runs over the points of the row without $u_j$), $u_k$ the third charge, and $C$ the projected colour factor of $\OO^{(2)}$ for that merge. The merges and kernels are those of FS Sec.~9 (the colour now with the general-$U$ placement). Merges not listed are zero.

@@THREE@@

%==========================================================================================
\section{Final result in colour-singlet form}\label{sec:final}
For hadron production ($\overline{MS}$, $d_\perp=2-\epsilon$, $n_f=0$), with $a=\frac{k^+}{\vee}$ and all $\Ls$ terms of the rows collected into the remainder (FS~(6.2): their sum is $\NN\log\frac{k^+}{\sqrt{\vee\Lambda}}(\JJ-\BB)$):
\begin{equation}
\boxed{\begin{aligned}
\frac{d^3N_h}{d^3k}={}&\int_{k^+/\vee}^1\frac{d\zeta}{\zeta^2}\,D_{h/g}(\zeta,\mu_F)\,\frac{d^3N_g}{d^3k}\Big(\frac{k^+}{\zeta},\frac{\mathbf k}{\zeta}\Big),\\
\frac{d^3N_g}{d^3k}={}&\LL[1]\big|_{\alpha_s(\mu_F)}\qquad\text{(Eq.~\eqref{eq:LO})}\\
&+\NN\int_{x,y,z,w,w'}e^{-ik\cdot(w'-w)}\,\PP(x,y)\Big[\frac12\log\frac{\vee}{\Lambda}\big(\JJ_s+\BB_s\big)+\log\frac{k^+}{\sqrt{\vee\Lambda}}\big(\JJ_s-\BB_s\big)\Big]\qquad\text{(evolution; remainder)}\\
&+\frac{\alpha_s}{2\pi}N_c\Big[\frac{67}{9}-\frac{\pi^2}{3}\Big]\LL[1]\qquad\text{(running coupling, CMW constant)}\\
&-\frac{\alpha_s}{2\pi}N_c\log\Big(1-\frac{k^+}{\vee}\Big)\log\Big(\frac{\vee}{k^+}-1\Big)\LL[1]
+\frac{\alpha_sN_c}{4\pi}\log\Big(1-\frac{k^+}{\vee}\Big)\LL\big[\Eh_X+\Eh_{X'}\big]\qquad\text{($V_\vee$ and Q1)}\\
&+\Big[-\frac{\alpha_s}{2\pi}\int_{a}^1\frac{d\xi}{\xi^2}P^{(r)}_{gg,+}(\xi)\,\LL_\xi\big[\ell_\Delta(\mu_F)\big]\new{-\Lv H_4}\Big]+\frac{d^3N_{\rm lin}}{d^3k}\qquad\text{(DGLAP, Eqs.~\eqref{eq:sDG}, \new{\eqref{eq:sH4}}, \eqref{eq:slin})}\\
&+\Big[\eqref{eq:s2a}+\eqref{eq:s2b}+\eqref{eq:s58}+\eqref{eq:sIII2}\Big]_{\Ls\to0}\qquad\text{(two-$\rho$ rows)}\\
&+\Big[\text{Eq.~\eqref{eq:three} for G1 II, III, V(1), the $\KK_{\rm IV}$ and $\KK_C$ families}\Big]_{\Ls\to0}\qquad\text{(three-$\rho$ rows, two-$\rho$ parts)}\\
&+\frac{d^3N_{3\rho,S}}{d^3k}\ \big(\propto\langle\rho\rho\rho\rangle\text{; zero for a Gaussian projectile}\big)+\frac{d^3N_{4\rho}}{d^3k}\ \text{(Q2)}\\
&+O(\alpha_s^3)+O\Big(\frac{\Lambda}{k^+}\Big).
\end{aligned}}
\label{eq:final}
\end{equation}
\begin{itemize}
\item $\JJ_s$, $\BB_s$: the singlet brackets EC~(3.4) (JIMWLK$\otimes$LO) and EC~(4.3) (BFKL$\otimes$LO). In them $x$, $y$ are the amplitude and conjugate charges and $w$, $w'$ the measured gluon (here: $x$, $x'$ and $y$, $y'$).
\item $\Ls\to0$: the rows are written without their $\Ls$ terms, which are all in the remainder $\log\frac{k^+}{\sqrt{\vee\Lambda}}(\JJ_s-\BB_s)$; it vanishes at mid rapidity. Equivalently, keep the $\Ls$ terms in the rows (and in Eq.~\eqref{eq:sLs}) and drop the remainder; not both.
\item Group I and Group III terms include their c.c.; Group II terms have none.
\item \new{$-\Lv H_4$: the $\log\frac{\vee}{k^+}$ at the lower end of the DGLAP integral is already in the evolution line ($\frac12\log\frac\vee\Lambda H_4$ in $\BB_s$, $\Ls H_4$ in the remainder). The DGLAP bracket is finite for $\frac{k^+}{\vee}\to0$. If the $\Ls$ terms are kept in the rows instead, the DGLAP bracket is $-\frac{\alpha_s}{2\pi}\int_a^1\dots-\frac12\log\frac\vee\Lambda H_4$.}
\item For option (ii) of Sec.~\ref{sec:log2} delete the $V_\vee$ and Q1 line.
\item LO is $O(\alpha_s)$; every other term is $O(\alpha_s^2)$. Nothing is expanded in $\frac{k^+}{\vee}$.
\end{itemize}

%==========================================================================================
\section{Checks}\label{sec:checks}
Scripts in \texttt{nlo\_gluon\_paper/finite\_singlet/} (they use \texttt{../singlet/trace.py}, \texttt{../finite\_v2/compact.py}, \texttt{../groups23/fixed23.py}, \texttt{../group1/user1.py}):
\begin{table}[H]\centering\small
\begin{tabular}{@{}p{0.66\textwidth}p{0.28\textwidth}@{}}\toprule
check & result\\\midrule
$\log^2$: $\Psi^A+\Psi^B$ (numerical $\xi$ integral) vs.\ the bracket of Row II part 1, $a=0.05\dots0.9$ (\texttt{log2\_origin.py}) & @@L2FIN@@ / @@L2POL@@\\
the same with the regular terms from $\xi=0$, vs.\ $\frac{\pi^2}{3}-\frac{67}{9}$ and no $\log(1-a)$ & @@L2UNR@@\\
the four closed-form integrals of Sec.~\ref{sec:origin} & @@L2PCS@@\\
general-$U$ operators (Groups II, III) vs.\ \texttt{fixed23.py}, general $U$ (\texttt{gsrc.py}) & @@GSMAX@@\\
the same vs.\ your operators for $U=U^T$ & @@GSSYM@@\\
Group I operators as written vs.\ \texttt{user1.py} (first principles), general $U$ & @@G1MAX@@\\
two-$\rho$ row colours vs.\ first-principles colour functions & @@FPTWO@@\\
two-$\rho$ parts by the pair rule vs.\ the reordering engine, general $U$ & @@ENG@@\\
every trace form vs.\ direct contraction (before / after combining differences) & @@TRMAX@@ / @@TRCMB@@\\
LO trace, Eq.~\eqref{eq:LO}, vs.\ the projected LO colour & @@LOCHK@@\\\bottomrule
\end{tabular}
\end{table}
All colour checks use random SU(3) adjoint Wilson lines with $U^T\ne U$ and three random configurations; the relative differences are with respect to the sum of the moduli of the single terms.
\end{document}
